\begin{table*}[!t]
\centering
\caption{Constants of the optional memory module.}
\label{tab:mmemconst}
\footnotesize
\renewcommand{\arraystretch}{1.15}
\begin{tabularx}{\textwidth}{@{}l r l c >{\raggedright\arraybackslash}X@{}}
\toprule
Symbol & Value & Unit & Basis & Meaning, justification, source\\
\midrule
$[\mathrm{Mem}]_\mathrm{tot}$ & \num{5000} & $\mathrm{molec/voxel}$ & E & Large pool so that the centroid of the mark is not shot-noise limited. \\
$k_\mathrm{w}$ & \num{4e-3} & $\mathrm{s^{-1}}$ & E & Writing rate, chosen so that less than $\approx\!40\,\%$ of the pool is written per wave (an unsaturated integrator). \\
$K_\mathrm{mem},\ n_\mathrm{mem}$ & \num{5000},\ 3 & $\mathrm{molec/voxel},\ -$ & C & PIP3 level of half-writing, above resting PIP3 and inside the range reached on the rising edge of a wave. \\
$k_\mathrm{e}$ & \num{5.56e-3} & $\mathrm{s^{-1}}$ & E & Erasure with $\tau=3$\,min, so that the mark persists through the back of a wave: cells keep their direction there for waves of the natural 6\,min period~\cite{skoge2014}. \\
$D_{\mathrm{Mem}_c},D_{\mathrm{Mem}_m}$ & \num{10},\ \num{0.01} & $\mu\mathrm{m^2/s}$ & E & Cytosolic carrier and nearly immobile mark. \\
\bottomrule
\end{tabularx}
\end{table*}
